\documentclass[aspectratio=169,11pt]{beamer}

\usetheme{Madrid}
\usecolortheme{default}
\usefonttheme{professionalfonts}
\setbeamertemplate{navigation symbols}{}
\setbeamertemplate{footline}[frame number]

\usepackage{amsmath, amssymb, booktabs, graphicx}

\newcommand{\E}{\mathbb{E}}
\newcommand{\Prb}{\mathbb{P}}

\title{Experiments and Field Experiments}
\subtitle{Empirical Methods --- Week 2}
\author{Teaching deck draft}
\date{}

\begin{document}

\begin{frame}
  \titlepage
\end{frame}

\begin{frame}{Week 2 question}
\begin{itemize}
    \item What does randomization identify?
    \item What remains hard after random assignment?
    \item How do compliance, spillovers, implementation, scale, and equilibrium change interpretation?
\end{itemize}

\medskip
Experiments are a design for building a counterfactual, not a ritual that ends the research argument.
\end{frame}

\begin{frame}{What randomization identifies}
\[
Z_i \perp \{Y_i(1),Y_i(0),X_i\}
\]

\[
\widehat{\tau}_{ATE}
=
\frac{1}{N_1}\sum_{i:Z_i=1}Y_i
-
\frac{1}{N_0}\sum_{i:Z_i=0}Y_i
\]

\begin{itemize}
    \item Assignment makes treated and control groups comparable in expectation.
    \item With full compliance and no interference, the mean contrast identifies the experimental ATE.
    \item Covariates can improve precision; randomization is the source of identification.
\end{itemize}
\end{frame}

\begin{frame}{Assignment is not receipt}
\begin{itemize}
    \item $Z_i$: randomized offer, encouragement, eligibility, information, or price.
    \item $D_i$: realized take-up or treatment receipt.
    \item If $D_i \neq Z_i$, the clean design object is the effect of assignment.
\end{itemize}

\medskip
Field experiments often randomize what researchers or policymakers can actually control: invitations, access, timing, caseworker scripts, prices, signals, or platform rules.
\end{frame}

\begin{frame}{ITT, take-up, and LATE}
\[
ITT_Y = \E[Y_i\mid Z_i=1] - \E[Y_i\mid Z_i=0]
\]

\[
ITT_D = \E[D_i\mid Z_i=1] - \E[D_i\mid Z_i=0]
\]

\[
\tau_{LATE}=\frac{ITT_Y}{ITT_D}
\]

\begin{itemize}
    \item ITT is often the policy object when the policy is an offer.
    \item The first stage measures how much assignment moves receipt.
    \item Wald/LATE is a complier effect under relevance, exclusion, and monotonicity.
\end{itemize}
\end{frame}

\begin{frame}{Duflo--Saez as design anchor}
\begin{itemize}
    \item Randomized encouragement to attend a retirement-plan benefits fair.
    \item Outcomes: attendance and retirement-plan participation.
    \item Design lessons:
    \begin{itemize}
        \item assignment differs from attendance;
        \item peer exposure can be part of the mechanism;
        \item ITT, first stage, and spillovers answer different questions.
    \end{itemize}
\end{itemize}
\end{frame}

\begin{frame}{Audit and correspondence experiments}
\begin{itemize}
    \item Randomize signals in real screening environments.
    \item Bertrand--Mullainathan: randomize names on resumes and measure callbacks.
    \item Strength: isolates employer response to a signal at initial screening.
    \item Limit: does not identify all hiring, bargaining, productivity, or equilibrium effects.
\end{itemize}
\end{frame}

\begin{frame}{Cluster randomization and spillovers}
\[
Y_i(\mathbf{Z}) = Y_i\left(Z_i, G_i(\mathbf{Z}_{-i})\right)
\]

\begin{itemize}
    \item Cluster randomization matches treatments delivered through firms, schools, offices, villages, or markets.
    \item Power depends heavily on the number of clusters and within-cluster correlation.
    \item Spillovers can contaminate a simple contrast or become the mechanism of interest.
    \item Saturation designs help separate direct effects from exposure effects.
\end{itemize}
\end{frame}

\begin{frame}{Firm and platform field experiments}
\small
\begin{tabular}{p{0.25\linewidth} p{0.31\linewidth} p{0.32\linewidth}}
\toprule
Paper role & What is randomized & Design lesson \\
\midrule
Bloom et al. & Management consulting in firms & Implementation and organizational response matter. \\
Pallais & Platform information / reputation & Digital markets can reveal information frictions. \\
Crepon et al. & Labor-market policy exposure & Direct gains can coexist with displacement. \\
\bottomrule
\end{tabular}
\end{frame}

\begin{frame}{External validity, scale, equilibrium}
\begin{itemize}
    \item Experimental population: who entered the study?
    \item Treatment intensity: is the pilot treatment the scalable policy?
    \item Institution: would delivery work outside this partner or platform?
    \item Market response: do prices, queues, vacancies, referrals, or competitors adjust?
    \item Welfare: are untreated workers or firms affected by the rollout?
\end{itemize}
\end{frame}

\begin{frame}{Robustness is design evidence}
\begin{itemize}
    \item Baseline balance and assignment protocol.
    \item Compliance, take-up, and first stage.
    \item Attrition and outcome measurement by assignment arm.
    \item Correct inference for individual, cluster, or market assignment.
    \item Spillover and exposure diagnostics.
    \item Power and minimum detectable effects for meaningful outcomes.
    \item Ethics, implementation logs, and reproducible randomization code.
\end{itemize}
\end{frame}

\begin{frame}{Research Lab design}
\begin{columns}[T,onlytextwidth]
\begin{column}{0.32\linewidth}
\textbf{Reproduce}
\begin{itemize}
    \item Synthetic Duflo--Saez-style encouragement data
    \item ITT, first stage, Wald/LATE
    \item Naive receipt contrast
\end{itemize}
\end{column}
\begin{column}{0.32\linewidth}
\textbf{Diagnose}
\begin{itemize}
    \item Balance
    \item Take-up
    \item Attrition
    \item Peer exposure
\end{itemize}
\end{column}
\begin{column}{0.32\linewidth}
\textbf{Transfer}
\begin{itemize}
    \item Labor-market saturation design
    \item Direct effects
    \item Spillovers
    \item Market-level interpretation
\end{itemize}
\end{column}
\end{columns}
\end{frame}

\begin{frame}{Takeaways}
\begin{itemize}
    \item Randomization solves assignment, not implementation or interpretation.
    \item ITT is the clean effect of assignment and often the policy-relevant object.
    \item LATE is local to compliers when encouragement moves receipt.
    \item Spillovers, clustered assignment, and equilibrium can change the estimand.
    \item Next week: DID and event studies build counterfactuals from time paths rather than random assignment.
\end{itemize}
\end{frame}

\end{document}
